% =======================================================
% 8. PREVENTIVO A FINIRE (PaF)
% =======================================================
\section{Preventivo a Finire (PaF)}
\label{sec:paf}

Il Preventivo a Finire viene aggiornato al termine di ogni sprint sulla base
degli scostamenti rilevati nei consuntivi. Il suo scopo è fornire una stima
aggiornata del costo complessivo del progetto, tenendo conto del costo già
sostenuto e del costo stimato per completare le attività rimanenti.

\subsection{Analisi scostamenti}

L'avanzamento della fase RTB\textsubscript{G} ha evidenziato uno scostamento
nella distribuzione delle ore tra i ruoli rispetto alla preventivazione iniziale.
In particolare, il carico associato alle attività di Amministrazione è risultato
superiore a quanto previsto, a causa della necessità di mantenere aggiornati in
modo continuativo gli strumenti di gestione, la documentazione di progetto, il
tracciamento delle attività e i consuntivi di sprint.

Parallelamente, alcune attività inizialmente associate a ruoli con tariffa oraria
maggiore, in particolare Analista\textsubscript{G} e Progettista\textsubscript{G},
si sono rivelate meno onerose del previsto nella fase corrente. Questo consente
di ribilanciare parte dell'impegno verso il ruolo di Amministratore senza
generare un incremento del costo complessivo stimato.

Lo scostamento rilevato riguarda quindi principalmente la distribuzione interna
delle ore tra ruoli e membri del gruppo, non il superamento del budget economico
preventivato.

\subsection{Ripianificazione}

A seguito degli scostamenti osservati, il gruppo ha aggiornato la pianificazione
delle attività rimanenti mantenendo invariato il vincolo economico complessivo
definito nel preventivo iniziale.

La ripianificazione prevede un ribilanciamento delle ore residue tra ruoli,
tenendo conto dell'effettivo andamento del progetto e delle attività ancora da
completare. In particolare, viene riconosciuto un maggiore peso alle attività di
Amministrazione, necessarie per il coordinamento operativo, la gestione degli
strumenti, l'aggiornamento della documentazione e il monitoraggio delle risorse.

Tale ribilanciamento non comporta un aumento del costo totale stimato, poiché
l'incremento relativo delle ore di Amministrazione viene compensato dalla
riduzione dell'impegno previsto su ruoli con tariffa oraria superiore. Il gruppo
continua pertanto a considerare il budget iniziale come vincolo non superabile e
mantiene il monitoraggio delle ore residue per ruolo e per componente al termine
di ogni sprint.

Eventuali ulteriori scostamenti verranno valutati nelle successive
retrospettive e recepiti nel Preventivo a Finire, aggiornando la stima delle
risorse residue senza modificare il perimetro economico complessivo del progetto.

\subsection{Quadro economico aggiornato}

\begin{table}[H]
    \centering
    \renewcommand{\arraystretch}{1.3}
    \begin{tabular}{l r}
        \toprule
        \textbf{Voce} & \textbf{Importo} \\
        \midrule
        Budget iniziale preventivato  & 11.665{,}00~\euro \\
        Costo sostenuto finora        & 2.295{,}58~\euro \\
        Costo stimato per completare  & 9.369{,}42~\euro \\
        \midrule
        \textbf{Nuovo totale stimato} & \textbf{11.665{,}00~\euro} \\
        \bottomrule
    \end{tabular}
    \caption{Quadro economico aggiornato del Preventivo a Finire dopo lo Sprint 7}
\end{table}